% ============================================================
% Descripcion textual de los casos de uso - App Turismo Comarapa
% Requiere en el preambulo del documento principal:
%   \usepackage{array}   % para columnas p{} y el modificador >{\bfseries}
%   \usepackage{float}   % para el especificador de posicion [H]
% ============================================================

% ------------------------------------------------------------
% MODULO PUBLICO (Actor: Visitante / Turista)
% ------------------------------------------------------------

\begin{table}[H]
\centering
\caption{Caso de uso CU-01 -- Consultar lugares turísticos}
\label{tab:cu01}
\renewcommand{\arraystretch}{1.3}
\begin{tabular}{|>{\bfseries}p{3.2cm}|p{10.3cm}|}
\hline
Actor & Visitante / Turista \\ \hline
Precondición & La aplicación está instalada y, preferentemente, el dispositivo cuenta con conexión a internet para obtener la información actualizada. \\ \hline
Flujo principal &
1. El usuario abre la aplicación y accede a la pantalla principal.\\
2. El usuario selecciona la categoría ``Lugares turísticos''.\\
3. El sistema consulta y muestra el listado de lugares activos registrados.\\
4. El usuario selecciona un lugar de la lista.\\
5. El sistema muestra el detalle del lugar (descripción, imágenes, ubicación y calificación promedio). \\ \hline
Flujos alternativos &
\textbullet\ Si no hay conexión a internet, el sistema muestra información de referencia (modo demo) o un mensaje de error con opción de reintentar.\\
\textbullet\ Si no existen lugares registrados, el sistema muestra un mensaje indicando que no hay resultados disponibles. \\ \hline
Postcondición & El usuario visualiza la información completa del lugar turístico seleccionado. \\ \hline
\end{tabular}
\end{table}

\begin{table}[H]
\centering
\caption{Caso de uso CU-02 -- Consultar actividades turísticas}
\label{tab:cu02}
\renewcommand{\arraystretch}{1.3}
\begin{tabular}{|>{\bfseries}p{3.2cm}|p{10.3cm}|}
\hline
Actor & Visitante / Turista \\ \hline
Precondición & La aplicación está instalada y, preferentemente, el dispositivo cuenta con conexión a internet. \\ \hline
Flujo principal &
1. El usuario accede a la pantalla principal.\\
2. El usuario selecciona la categoría ``Actividades''.\\
3. El sistema consulta y muestra el listado de actividades activas registradas.\\
4. El usuario selecciona una actividad de la lista.\\
5. El sistema muestra el detalle de la actividad (descripción, imágenes y ubicación). \\ \hline
Flujos alternativos &
\textbullet\ Si no hay conexión, el sistema muestra un mensaje de error con opción de reintentar.\\
\textbullet\ Si no existen actividades registradas, se muestra un mensaje de ``sin resultados''. \\ \hline
Postcondición & El usuario visualiza la información completa de la actividad seleccionada. \\ \hline
\end{tabular}
\end{table}

\begin{table}[H]
\centering
\caption{Caso de uso CU-03 -- Consultar eventos y ferias}
\label{tab:cu03}
\renewcommand{\arraystretch}{1.3}
\begin{tabular}{|>{\bfseries}p{3.2cm}|p{10.3cm}|}
\hline
Actor & Visitante / Turista \\ \hline
Precondición & La aplicación está instalada y, preferentemente, el dispositivo cuenta con conexión a internet. \\ \hline
Flujo principal &
1. El usuario accede a la pantalla principal.\\
2. El usuario selecciona la categoría ``Eventos'' (ferias, fiestas patronales, actividades culturales).\\
3. El sistema consulta y muestra el listado de eventos activos, incluyendo fecha y periodicidad (anual o único).\\
4. El usuario selecciona un evento de la lista.\\
5. El sistema muestra el detalle del evento. \\ \hline
Flujos alternativos &
\textbullet\ Si un evento fue marcado como inactivo por el administrador, no se muestra en el listado público.\\
\textbullet\ Si no existen eventos próximos, se muestra un mensaje de ``sin resultados''. \\ \hline
Postcondición & El usuario visualiza la información completa del evento o feria seleccionado. \\ \hline
\end{tabular}
\end{table}

\begin{table}[H]
\centering
\caption{Caso de uso CU-04 -- Consultar hospedaje}
\label{tab:cu04}
\renewcommand{\arraystretch}{1.3}
\begin{tabular}{|>{\bfseries}p{3.2cm}|p{10.3cm}|}
\hline
Actor & Visitante / Turista \\ \hline
Precondición & La aplicación está instalada y, preferentemente, el dispositivo cuenta con conexión a internet. \\ \hline
Flujo principal &
1. El usuario accede a la pantalla principal.\\
2. El usuario selecciona la categoría ``Hoteles''.\\
3. El sistema consulta y muestra el listado de establecimientos de hospedaje activos.\\
4. El usuario selecciona un hotel de la lista.\\
5. El sistema muestra el detalle del hotel (descripción, precio referencial, imágenes y ubicación). \\ \hline
Flujos alternativos &
\textbullet\ Si no hay conexión, se muestra un mensaje de error con opción de reintentar.\\
\textbullet\ Si no existen hoteles registrados, se muestra un mensaje de ``sin resultados''. \\ \hline
Postcondición & El usuario visualiza la información completa del establecimiento de hospedaje seleccionado. \\ \hline
\end{tabular}
\end{table}

\begin{table}[H]
\centering
\caption{Caso de uso CU-05 -- Consultar restaurantes y gastronomía}
\label{tab:cu05}
\renewcommand{\arraystretch}{1.3}
\begin{tabular}{|>{\bfseries}p{3.2cm}|p{10.3cm}|}
\hline
Actor & Visitante / Turista \\ \hline
Precondición & La aplicación está instalada y, preferentemente, el dispositivo cuenta con conexión a internet. \\ \hline
Flujo principal &
1. El usuario accede a la pantalla principal.\\
2. El usuario selecciona la categoría ``Restaurantes'' o ``Gastronomía''.\\
3. El sistema consulta y muestra el listado de restaurantes o platillos típicos activos.\\
4. El usuario selecciona un restaurante o platillo de la lista.\\
5. El sistema muestra el detalle correspondiente; en el caso de un restaurante, incluye los platillos típicos asociados a él. \\ \hline
Flujos alternativos &
\textbullet\ Si no hay conexión, se muestra un mensaje de error con opción de reintentar.\\
\textbullet\ Si un restaurante no tiene platillos registrados, la sección de gastronomía asociada se muestra vacía. \\ \hline
Postcondición & El usuario visualiza la información completa del restaurante o platillo seleccionado. \\ \hline
\end{tabular}
\end{table}

\begin{table}[H]
\centering
\caption{Caso de uso CU-06 -- Buscar contenido turístico}
\label{tab:cu06}
\renewcommand{\arraystretch}{1.3}
\begin{tabular}{|>{\bfseries}p{3.2cm}|p{10.3cm}|}
\hline
Actor & Visitante / Turista \\ \hline
Precondición & Existen registros activos en al menos una de las seis categorías turísticas. \\ \hline
Flujo principal &
1. El usuario accede a la pantalla de búsqueda.\\
2. El usuario ingresa un texto en el campo de búsqueda.\\
3. El sistema consulta simultáneamente lugares, actividades, eventos, hoteles, restaurantes y gastronomía filtrando por coincidencia de texto.\\
4. El sistema normaliza y combina los resultados en una sola lista.\\
5. El usuario selecciona un resultado.\\
6. El sistema navega al detalle correspondiente según la categoría del resultado. \\ \hline
Flujos alternativos &
\textbullet\ Si no existen coincidencias, el sistema muestra el mensaje ``sin resultados''.\\
\textbullet\ Si el campo de búsqueda está vacío, el sistema no ejecuta la búsqueda y mantiene el estado inicial. \\ \hline
Postcondición & El usuario visualiza el detalle del elemento turístico encontrado. \\ \hline
\end{tabular}
\end{table}

\begin{table}[H]
\centering
\caption{Caso de uso CU-07 -- Visualizar mapa interactivo}
\label{tab:cu07}
\renewcommand{\arraystretch}{1.3}
\begin{tabular}{|>{\bfseries}p{3.2cm}|p{10.3cm}|}
\hline
Actor & Visitante / Turista \\ \hline
Precondición & El dispositivo cuenta con conexión a internet para descargar las teselas del mapa. \\ \hline
Flujo principal &
1. El usuario abre la sección de mapa.\\
2. El sistema descarga y muestra el mapa (OpenStreetMap) centrado en el municipio de Comarapa.\\
3. El sistema coloca marcadores diferenciados por categoría (lugares, hoteles, restaurantes, eventos).\\
4. El usuario filtra los marcadores visibles mediante chips de categoría.\\
5. El usuario toca un marcador y el sistema muestra una tarjeta de previsualización con los datos básicos del punto. \\ \hline
Flujos alternativos &
\textbullet\ Si falla la descarga de las teselas por falta de conexión, se muestra un indicador de carga o un estado de error.\\
\textbullet\ El usuario puede, desde esta pantalla, extender el caso de uso hacia CU-08 (ubicar posición actual) o CU-09 (calcular ruta hacia destino). \\ \hline
Postcondición & El usuario visualiza en el mapa la ubicación geográfica de los puntos turísticos del municipio. \\ \hline
\end{tabular}
\end{table}

\begin{table}[H]
\centering
\caption{Caso de uso CU-08 -- Ubicar posición actual (GPS) \textit{(extiende a CU-07)}}
\label{tab:cu08}
\renewcommand{\arraystretch}{1.3}
\begin{tabular}{|>{\bfseries}p{3.2cm}|p{10.3cm}|}
\hline
Actor & Visitante / Turista \\ \hline
Precondición & El usuario otorgó permiso de ubicación a la aplicación y el GPS del dispositivo está activo. \\ \hline
Flujo principal &
1. El usuario presiona el botón flotante ``Mi ubicación'' dentro del mapa.\\
2. El sistema verifica el permiso de ubicación concedido.\\
3. El sistema obtiene las coordenadas actuales del dispositivo mediante el servicio de geolocalización.\\
4. El sistema centra el mapa en la posición obtenida y coloca un marcador de usuario. \\ \hline
Flujos alternativos &
\textbullet\ Si el permiso de ubicación es denegado, el sistema muestra un mensaje solicitando habilitarlo desde la configuración del dispositivo.\\
\textbullet\ Si el GPS está desactivado, el sistema solicita al usuario activarlo. \\ \hline
Postcondición & El mapa queda centrado en la ubicación actual del usuario. \\ \hline
\end{tabular}
\end{table}

\begin{table}[H]
\centering
\caption{Caso de uso CU-09 -- Calcular ruta hacia destino \textit{(extiende a CU-07)}}
\label{tab:cu09}
\renewcommand{\arraystretch}{1.3}
\begin{tabular}{|>{\bfseries}p{3.2cm}|p{10.3cm}|}
\hline
Actor & Visitante / Turista \\ \hline
Precondición & Se cuenta con la ubicación actual del usuario (CU-08) y existe conexión a internet para consultar el servicio de enrutamiento. \\ \hline
Flujo principal &
1. El usuario selecciona la opción ``Cómo llegar'' desde el detalle de un lugar o desde el mapa.\\
2. El sistema toma la ubicación actual del usuario y las coordenadas del destino.\\
3. El sistema envía la solicitud al servicio de enrutamiento (OSRM).\\
4. El sistema recibe y dibuja la ruta sobre el mapa, junto con la distancia y el tiempo estimado. \\ \hline
Flujos alternativos &
\textbullet\ Si el servicio de enrutamiento no responde o no hay conexión a internet, el sistema muestra un mensaje indicando que no fue posible calcular la ruta. \\ \hline
Postcondición & El usuario visualiza la ruta trazada entre su ubicación y el destino turístico seleccionado. \\ \hline
\end{tabular}
\end{table}

\begin{table}[H]
\centering
\caption{Caso de uso CU-10 -- Consultar el clima}
\label{tab:cu10}
\renewcommand{\arraystretch}{1.3}
\begin{tabular}{|>{\bfseries}p{3.2cm}|p{10.3cm}|}
\hline
Actor & Visitante / Turista \\ \hline
Precondición & El dispositivo cuenta con conexión a internet para consultar el servicio meteorológico. \\ \hline
Flujo principal &
1. El sistema, al cargar la pantalla principal, solicita el clima actual del municipio de Comarapa al servicio meteorológico.\\
2. El sistema recibe la información (temperatura y condición climática).\\
3. El sistema muestra el resumen del clima en la pantalla principal. \\ \hline
Flujos alternativos &
\textbullet\ Si el servicio meteorológico no responde, el sistema oculta o muestra vacío el bloque de clima sin interrumpir el resto de la pantalla, permitiendo reintentar. \\ \hline
Postcondición & El usuario visualiza la condición climática actual del municipio. \\ \hline
\end{tabular}
\end{table}

\begin{table}[H]
\centering
\caption{Caso de uso CU-11 -- Registrar reseña y calificación}
\label{tab:cu11}
\renewcommand{\arraystretch}{1.3}
\begin{tabular}{|>{\bfseries}p{3.2cm}|p{10.3cm}|}
\hline
Actor & Visitante / Turista \\ \hline
Precondición & El usuario se encuentra en la pantalla de detalle de un lugar, hotel, restaurante, actividad, gastronomía o evento. No requiere una cuenta registrada. \\ \hline
Flujo principal &
1. El usuario presiona la opción ``Escribir reseña'' dentro de la sección de reseñas.\\
2. El usuario ingresa su nombre, selecciona una calificación de 1 a 5 estrellas y, opcionalmente, escribe un comentario.\\
3. El sistema valida los datos ingresados (nombre y comentario no exceden la longitud permitida).\\
4. El sistema registra la reseña asociada a la entidad turística correspondiente.\\
5. El sistema actualiza el promedio de calificación y el listado de reseñas visibles. \\ \hline
Flujos alternativos &
\textbullet\ Si el comentario supera los 500 caracteres, el sistema muestra un mensaje de validación y no permite guardar.\\
\textbullet\ Si el usuario cancela el formulario, la reseña no se registra. \\ \hline
Postcondición & La reseña queda registrada y visible junto con el promedio actualizado de calificación de la entidad. \\ \hline
\end{tabular}
\end{table}

\begin{table}[H]
\centering
\caption{Caso de uso CU-12 -- Consultar reseñas y promedio}
\label{tab:cu12}
\renewcommand{\arraystretch}{1.3}
\begin{tabular}{|>{\bfseries}p{3.2cm}|p{10.3cm}|}
\hline
Actor & Visitante / Turista \\ \hline
Precondición & La entidad turística consultada existe en el sistema. \\ \hline
Flujo principal &
1. El usuario abre la pantalla de detalle de una entidad turística.\\
2. El sistema consulta las reseñas asociadas a dicha entidad.\\
3. El sistema calcula el promedio de calificación y el total de reseñas.\\
4. El sistema muestra el listado de reseñas (autor, calificación, comentario y fecha) junto con el promedio. \\ \hline
Flujos alternativos &
\textbullet\ Si la entidad aún no tiene reseñas, el sistema muestra el mensaje ``Aún no hay reseñas'' y un promedio de 0. \\ \hline
Postcondición & El usuario visualiza la opinión y calificación de otros visitantes sobre la entidad turística. \\ \hline
\end{tabular}
\end{table}

% ------------------------------------------------------------
% MODULO ADMINISTRATIVO (Actores: Editor y Administrador)
% ------------------------------------------------------------

\begin{table}[H]
\centering
\caption{Caso de uso CU-13 -- Iniciar sesión}
\label{tab:cu13}
\renewcommand{\arraystretch}{1.3}
\begin{tabular}{|>{\bfseries}p{3.2cm}|p{10.3cm}|}
\hline
Actor & Editor, Administrador \\ \hline
Precondición & El usuario cuenta con una cuenta registrada en el sistema de autenticación y no tiene una sesión activa. \\ \hline
Flujo principal &
1. El usuario accede a la pantalla de inicio de sesión del panel administrativo.\\
2. El usuario ingresa su correo electrónico y contraseña.\\
3. El sistema valida el formato de los campos ingresados.\\
4. El sistema envía las credenciales al servicio de autenticación.\\
5. El sistema verifica las credenciales y, de ser correctas, crea la sesión y da acceso al panel administrativo según el rol del usuario. \\ \hline
Flujos alternativos &
\textbullet\ Si las credenciales son incorrectas, el sistema muestra el mensaje ``Correo o contraseña incorrectos''.\\
\textbullet\ Si el correo no ha sido confirmado, ocurre un error de red, o se excede el número de intentos, el sistema muestra el mensaje específico correspondiente y permite reintentar. \\ \hline
Postcondición & El usuario queda autenticado y accede al panel de administración correspondiente a su rol. \\ \hline
\end{tabular}
\end{table}

\begin{table}[H]
\centering
\caption{Caso de uso CU-14 -- Registrarse como editor}
\label{tab:cu14}
\renewcommand{\arraystretch}{1.3}
\begin{tabular}{|>{\bfseries}p{3.2cm}|p{10.3cm}|}
\hline
Actor & Visitante (usuario que solicita una cuenta) \\ \hline
Precondición & El usuario cuenta con un correo electrónico válido no registrado previamente en el sistema. \\ \hline
Flujo principal &
1. El usuario selecciona la opción de registro desde la pantalla de inicio de sesión.\\
2. El usuario ingresa correo electrónico y contraseña (mínimo 6 caracteres).\\
3. El sistema valida el formulario.\\
4. El sistema envía la solicitud de registro al servicio de autenticación.\\
5. El sistema crea la cuenta y, mediante un disparador de la base de datos, le asigna automáticamente el rol ``editor''. \\ \hline
Flujos alternativos &
\textbullet\ Si el correo ya se encuentra registrado, la contraseña es débil, o el registro está deshabilitado, el sistema muestra el mensaje de error correspondiente. \\ \hline
Postcondición & Se crea una nueva cuenta de usuario con rol ``editor'', pendiente de iniciar sesión (CU-13). \\ \hline
\end{tabular}
\end{table}

\begin{table}[H]
\centering
\caption{Caso de uso CU-15 -- Cerrar sesión}
\label{tab:cu15}
\renewcommand{\arraystretch}{1.3}
\begin{tabular}{|>{\bfseries}p{3.2cm}|p{10.3cm}|}
\hline
Actor & Editor, Administrador \\ \hline
Precondición & El usuario tiene una sesión activa en el panel administrativo. \\ \hline
Flujo principal &
1. El usuario selecciona la opción ``Cerrar sesión'' dentro del panel administrativo.\\
2. El sistema finaliza la sesión activa en el servicio de autenticación.\\
3. El sistema redirige al usuario a la pantalla de inicio de sesión. \\ \hline
Flujos alternativos &
\textbullet\ No se identifican flujos alternativos relevantes para esta operación. \\ \hline
Postcondición & La sesión del usuario queda cerrada; se requiere volver a autenticarse para acceder nuevamente al panel administrativo. \\ \hline
\end{tabular}
\end{table}

\begin{table}[H]
\centering
\caption{Caso de uso CU-16 -- Gestionar lugares turísticos}
\label{tab:cu16}
\renewcommand{\arraystretch}{1.3}
\begin{tabular}{|>{\bfseries}p{3.2cm}|p{10.3cm}|}
\hline
Actor & Editor, Administrador \\ \hline
Precondición & El usuario inició sesión (CU-13) con rol editor o administrador. \\ \hline
Flujo principal &
1. El usuario accede a la sección ``Lugares'' del panel administrativo.\\
2. El sistema muestra el listado de lugares registrados, con opciones de búsqueda y filtrado por estado.\\
3. El usuario elige crear un nuevo lugar, o editar/eliminar uno existente.\\
4. El usuario completa el formulario (nombre, categoría, descripción, imágenes y ubicación en el mapa).\\
5. El sistema valida los datos ingresados y los guarda en la base de datos.\\
6. El sistema actualiza el listado administrativo y refleja los cambios en el módulo público. \\ \hline
Flujos alternativos &
\textbullet\ Si el usuario cancela la operación, no se guardan los cambios.\\
\textbullet\ Si un campo obligatorio es inválido, el sistema muestra el error de validación correspondiente.\\
\textbullet\ El usuario puede activar o desactivar la visibilidad pública del lugar sin eliminarlo definitivamente. \\ \hline
Postcondición & El lugar turístico queda creado, actualizado o eliminado, reflejándose en el módulo público según su estado. \\ \hline
\end{tabular}
\end{table}

\begin{table}[H]
\centering
\caption{Caso de uso CU-17 -- Gestionar hoteles}
\label{tab:cu17}
\renewcommand{\arraystretch}{1.3}
\begin{tabular}{|>{\bfseries}p{3.2cm}|p{10.3cm}|}
\hline
Actor & Editor, Administrador \\ \hline
Precondición & El usuario inició sesión (CU-13) con rol editor o administrador. \\ \hline
Flujo principal &
1. El usuario accede a la sección ``Hoteles'' del panel administrativo.\\
2. El sistema muestra el listado de establecimientos de hospedaje registrados.\\
3. El usuario elige crear un nuevo hotel, o editar/eliminar uno existente.\\
4. El usuario completa el formulario (nombre, descripción, precio referencial, imágenes y ubicación).\\
5. El sistema valida los datos ingresados y los guarda en la base de datos.\\
6. El sistema actualiza el listado administrativo y refleja los cambios en el módulo público. \\ \hline
Flujos alternativos &
\textbullet\ Si el usuario cancela la operación, no se guardan los cambios.\\
\textbullet\ Si un campo obligatorio es inválido, el sistema muestra el error de validación correspondiente.\\
\textbullet\ El usuario puede activar o desactivar la visibilidad pública del hotel sin eliminarlo definitivamente. \\ \hline
Postcondición & El hotel queda creado, actualizado o eliminado, reflejándose en el módulo público según su estado. \\ \hline
\end{tabular}
\end{table}

\begin{table}[H]
\centering
\caption{Caso de uso CU-18 -- Gestionar restaurantes}
\label{tab:cu18}
\renewcommand{\arraystretch}{1.3}
\begin{tabular}{|>{\bfseries}p{3.2cm}|p{10.3cm}|}
\hline
Actor & Editor, Administrador \\ \hline
Precondición & El usuario inició sesión (CU-13) con rol editor o administrador. \\ \hline
Flujo principal &
1. El usuario accede a la sección ``Restaurantes'' del panel administrativo.\\
2. El sistema muestra el listado de restaurantes registrados.\\
3. El usuario elige crear un nuevo restaurante, o editar/eliminar uno existente.\\
4. El usuario completa el formulario (nombre, descripción, imágenes y ubicación).\\
5. El sistema valida los datos ingresados y los guarda en la base de datos.\\
6. El sistema actualiza el listado administrativo y refleja los cambios en el módulo público. \\ \hline
Flujos alternativos &
\textbullet\ Si el usuario cancela la operación, no se guardan los cambios.\\
\textbullet\ Si un campo obligatorio es inválido, el sistema muestra el error de validación correspondiente.\\
\textbullet\ El usuario puede activar o desactivar la visibilidad pública del restaurante sin eliminarlo definitivamente. \\ \hline
Postcondición & El restaurante queda creado, actualizado o eliminado, reflejándose en el módulo público según su estado. \\ \hline
\end{tabular}
\end{table}

\begin{table}[H]
\centering
\caption{Caso de uso CU-19 -- Gestionar eventos y ferias}
\label{tab:cu19}
\renewcommand{\arraystretch}{1.3}
\begin{tabular}{|>{\bfseries}p{3.2cm}|p{10.3cm}|}
\hline
Actor & Editor, Administrador \\ \hline
Precondición & El usuario inició sesión (CU-13) con rol editor o administrador. \\ \hline
Flujo principal &
1. El usuario accede a la sección ``Eventos'' del panel administrativo.\\
2. El sistema muestra el listado de eventos y ferias registrados.\\
3. El usuario elige crear un nuevo evento, o editar/eliminar uno existente.\\
4. El usuario completa el formulario (nombre, categoría, fecha de inicio, periodicidad, descripción e imágenes).\\
5. El sistema valida los datos ingresados y los guarda en la base de datos.\\
6. El sistema actualiza el listado administrativo y refleja los cambios en el módulo público. \\ \hline
Flujos alternativos &
\textbullet\ Si el usuario cancela la operación, no se guardan los cambios.\\
\textbullet\ Si un campo obligatorio es inválido, el sistema muestra el error de validación correspondiente.\\
\textbullet\ El usuario puede activar o desactivar la visibilidad pública del evento sin eliminarlo definitivamente. \\ \hline
Postcondición & El evento o feria queda creado, actualizado o eliminado, reflejándose en el módulo público según su estado. \\ \hline
\end{tabular}
\end{table}

\begin{table}[H]
\centering
\caption{Caso de uso CU-20 -- Gestionar actividades}
\label{tab:cu20}
\renewcommand{\arraystretch}{1.3}
\begin{tabular}{|>{\bfseries}p{3.2cm}|p{10.3cm}|}
\hline
Actor & Editor, Administrador \\ \hline
Precondición & El usuario inició sesión (CU-13) con rol editor o administrador. \\ \hline
Flujo principal &
1. El usuario accede a la sección ``Actividades'' del panel administrativo.\\
2. El sistema muestra el listado de actividades turísticas registradas.\\
3. El usuario elige crear una nueva actividad, o editar/eliminar una existente.\\
4. El usuario completa el formulario (nombre, descripción, imágenes y ubicación).\\
5. El sistema valida los datos ingresados y los guarda en la base de datos.\\
6. El sistema actualiza el listado administrativo y refleja los cambios en el módulo público. \\ \hline
Flujos alternativos &
\textbullet\ Si el usuario cancela la operación, no se guardan los cambios.\\
\textbullet\ Si un campo obligatorio es inválido, el sistema muestra el error de validación correspondiente.\\
\textbullet\ El usuario puede activar o desactivar la visibilidad pública de la actividad sin eliminarla definitivamente. \\ \hline
Postcondición & La actividad turística queda creada, actualizada o eliminada, reflejándose en el módulo público según su estado. \\ \hline
\end{tabular}
\end{table}

\begin{table}[H]
\centering
\caption{Caso de uso CU-21 -- Gestionar gastronomía}
\label{tab:cu21}
\renewcommand{\arraystretch}{1.3}
\begin{tabular}{|>{\bfseries}p{3.2cm}|p{10.3cm}|}
\hline
Actor & Editor, Administrador \\ \hline
Precondición & El usuario inició sesión (CU-13) con rol editor o administrador. \\ \hline
Flujo principal &
1. El usuario accede a la sección ``Gastronomía'' del panel administrativo.\\
2. El sistema muestra el listado de platillos típicos registrados, opcionalmente asociados a un restaurante.\\
3. El usuario elige crear un nuevo platillo, o editar/eliminar uno existente.\\
4. El usuario completa el formulario (nombre, descripción, restaurante asociado e imágenes).\\
5. El sistema valida los datos ingresados y los guarda en la base de datos.\\
6. El sistema actualiza el listado administrativo y refleja los cambios en el módulo público. \\ \hline
Flujos alternativos &
\textbullet\ Si el usuario cancela la operación, no se guardan los cambios.\\
\textbullet\ Si un campo obligatorio es inválido, el sistema muestra el error de validación correspondiente.\\
\textbullet\ El usuario puede activar o desactivar la visibilidad pública del platillo sin eliminarlo definitivamente. \\ \hline
Postcondición & El registro de gastronomía queda creado, actualizado o eliminado, reflejándose en el módulo público según su estado. \\ \hline
\end{tabular}
\end{table}

\begin{table}[H]
\centering
\caption{Caso de uso CU-22 -- Seleccionar ubicación en el mapa \textit{(incluido en CU-16 a CU-21)}}
\label{tab:cu22}
\renewcommand{\arraystretch}{1.3}
\begin{tabular}{|>{\bfseries}p{3.2cm}|p{10.3cm}|}
\hline
Actor & Editor, Administrador \\ \hline
Precondición & El usuario se encuentra creando o editando un registro turístico que requiere coordenadas geográficas. \\ \hline
Flujo principal &
1. El usuario presiona el campo de ubicación dentro del formulario.\\
2. El sistema abre un mapa interactivo a pantalla completa, centrado en el municipio o en la ubicación previamente registrada.\\
3. El usuario desplaza el mapa o toca un punto para mover el marcador.\\
4. El usuario puede centrar el mapa en su posición GPS actual como referencia.\\
5. El usuario confirma la ubicación seleccionada.\\
6. El sistema retorna las coordenadas (latitud y longitud) al formulario de origen. \\ \hline
Flujos alternativos &
\textbullet\ Si el usuario cancela la selección, se conserva la ubicación previa o el campo queda vacío. \\ \hline
Postcondición & El registro turístico queda asociado a una coordenada geográfica válida. \\ \hline
\end{tabular}
\end{table}

\begin{table}[H]
\centering
\caption{Caso de uso CU-23 -- Subir imágenes \textit{(incluido en CU-16 a CU-21)}}
\label{tab:cu23}
\renewcommand{\arraystretch}{1.3}
\begin{tabular}{|>{\bfseries}p{3.2cm}|p{10.3cm}|}
\hline
Actor & Editor, Administrador \\ \hline
Precondición & El usuario se encuentra creando o editando un registro turístico y ha concedido permisos de cámara o galería. \\ \hline
Flujo principal &
1. El usuario presiona el campo de imágenes dentro del formulario.\\
2. El sistema muestra las opciones ``Cámara'' o ``Galería''.\\
3. El usuario captura o selecciona una fotografía.\\
4. El sistema sube el archivo al almacenamiento en la nube (bucket de imágenes).\\
5. El sistema obtiene la URL pública de la imagen y la agrega a la lista de imágenes del registro.\\
6. El sistema muestra la imagen agregada; la primera imagen de la lista se utiliza como portada en las pantallas públicas. \\ \hline
Flujos alternativos &
\textbullet\ Si el usuario cancela la selección, no se sube ninguna imagen.\\
\textbullet\ Si la subida falla por falta de conexión, el sistema muestra un mensaje de error y permite reintentar. \\ \hline
Postcondición & El registro turístico queda con una o más imágenes asociadas y disponibles públicamente. \\ \hline
\end{tabular}
\end{table}

\begin{table}[H]
\centering
\caption{Caso de uso CU-24 -- Gestionar usuarios (solo administrador)}
\label{tab:cu24}
\renewcommand{\arraystretch}{1.3}
\begin{tabular}{|>{\bfseries}p{3.2cm}|p{10.3cm}|}
\hline
Actor & Administrador \\ \hline
Precondición & El usuario inició sesión (CU-13) con rol ``administrador''. \\ \hline
Flujo principal &
1. El administrador accede a la sección ``Usuarios'' del panel administrativo.\\
2. El sistema muestra el listado de usuarios registrados con su rol y estado (activo/inactivo).\\
3. El administrador filtra el listado por rol (todos, administradores o editores).\\
4. El administrador selecciona un usuario y modifica su rol (editor/administrador) o su estado (activo/inactivo).\\
5. El sistema actualiza el registro del usuario en la base de datos. \\ \hline
Flujos alternativos &
\textbullet\ Si la conexión falla, el sistema muestra un listado de datos de referencia (demo) para mantener la fidelidad visual.\\
\textbullet\ El administrador puede, adicionalmente, registrar una nueva cuenta de usuario desde este mismo panel. \\ \hline
Postcondición & El rol o el estado del usuario queda actualizado en el sistema. \\ \hline
\end{tabular}
\end{table}
